\originalsection{一次完全积分怎样恢复初值解}
以一维谐振子作完整计算, $H=p^2/(2m)+m\omega^2q^2/2$, 取参数$E>0$, 振幅$A=\sqrt{2E/(m\omega^2)}$. 在$p>0$的局部分支上取
\[
W(q,E)=\int_0^q m\omega\sqrt{A^2-x^2}\dd x
=\frac{m\omega}{2}\left[q\sqrt{A^2-q^2}+A^2\arcsin(q/A)\right].
\]
令$S=W-Et$, 则$S_q=m\omega\sqrt{A^2-q^2}$, 代入$H$恰为$E$, 故HJ方程成立. 又$W_E=\omega^{-1}\arcsin(q/A)$, 所以$\beta=S_E=\omega^{-1}\arcsin(q/A)-t$是常数. 反解给$q=A\sin[\omega(t+\beta)]$, 再由$p=S_q$得当前分支上的$p=m\omega A\cos[\omega(t+\beta)]>0$. 沿轨道换到负动量分支后, 连续延拓相位, 就恢复整个周期解.

给定$(q_0,p_0)$和时刻$t_0$, 按$p_0$的符号选用对应动量分支, 再取$E=p_0^2/(2m)+m\omega^2q_0^2/2$, 再取初相位$\delta_0=\operatorname{atan2}(m\omega q_0,p_0)$, 则$\beta=\delta_0/\omega-t_0$模周期. 因而两个积分常数来自两个初值; 选取一个固定的$E$后不能还宣称描述所有初值. 在非转向点$S_{qE}=m/p\ne0$, 完全积分满足局部可逆条件. $p=0$处该分支表达不光滑, 须改用相邻生成函数坐标片; 原运动解本身并不奇异. $E=0$为固定点, 不属于上述非零振幅的角坐标.

沿真实路径,$\dd S/\dd t=S_q\dot q+S_t=p\dot q-H=L$. 所以$S$在路径上的差是作用量, 不是凭名字称“主函数”. 对不显含时间的$W$, 有$\dd W/\dd t=p\dot q$; $W$和$S$承担的时间项不同.

\originalsection{摆的转动作用量与分离轨道}
摆在$E>2mg\ell$时可连续转动. 取正方向动量, 一个角周期的作用量为
\[
I_{\rm rot}(E)=\frac1{2\pi}\int_0^{2\pi}
\sqrt{2m\ell^2[E-mg\ell(1-\cos\theta)]}\dd\theta.
\]
它与振荡作用量分别对应圆柱相空间中的两种回路. 临界能量$E_s=2mg\ell$上$p=2m\ell^2\sqrt{g/\ell}\,|\cos(\theta/2)|$. 因而从旋转侧趋近的$I_{\rm rot}\to4m\ell^2\sqrt{g/\ell}/\pi$, 从振荡侧趋近的$I_{\rm lib}\to8m\ell^2\sqrt{g/\ell}/\pi$. 后者回路还包含负动量支, 面积是正向旋转回路积分的两倍. 作用量不能跨越分离轨道不加说明地沿用同一个数值标签. 两侧周期均因顶部附近$\dot\theta\to0$而发散, 作用角坐标在这里退化; 这也预告绝热不变量在慢参数扫过分离轨道时需要特殊处理.
